% !PRINTED 169
\weberpage{169}
%<<CONT
ferner das System $\mathfrak{S}$ Formen 0\textsuperscript{ten} Grades, d.\,h. von Null verschiedene Constanten enthält, so ist, da wir für $F_1$ eine solche Constante nehmen können, unser Satz selbstverständlich, da, wenn $A_1 = F : F_1$ gesetzt wird, $F = A_1F_1$ ist. In dieser Form ist aber der Satz inhaltlos. Wenn wir aber von $\mathfrak{S}$ alle constanten Formen ausschliessen, so bleibt ein System $\mathfrak{S}'$, das keine Formen 0\textsuperscript{ten} Grades mehr enthält, für das unser Satz gleichfalls gilt. Die Basis $F_1, F_2, \ldots F_\mu$ enthält dann gleichfalls keine Formen 0\textsuperscript{ten} Grades, und wir können daher den Satz I. auch so ausdrücken:

\gesperrt{II. Alle Formen des Systemes $\mathfrak{S}$ von positivem Grade lassen sich durch eine Basis $F_1, F_2, \ldots F_\mu$, deren Elemente von positivem Grade sind, in der Form ausdrücken:}
\weq{1}{F = \Phi_1F_1 + \Phi_2F_2 + \cdots + \Phi_\mu F_\mu.}
\gesperrt{Die Grade von $\Phi_1, \Phi_2, \ldots \Phi_\mu$ sind dann niedriger als der Grad von $F$.}

Die wichtigsten Anwendungen findet dieses Theorem bei Untersuchungen über die Möglichkeit, alle Formen eines gewissen Systemes $\mathfrak{S}$ als \gesperrt{ganze rationale Functionen einer endlichen Anzahl} unter ihnen darzustellen. Man nennt ein solches Formensystem ein \gesperrt{endliches} (nicht in dem Sinne, dass es nur aus einer endlichen Zahl von Formen besteht). Es gilt der folgende Satz:

\gesperrt{III. Wenn sich die Co\"efficienten $\Phi_1, \Phi_2, \ldots \Phi_\mu$ in der Darstellung (1) des Theorems II. für jede Form $F$ in $\mathfrak{S}$ so wählen lassen, dass sie, wenn sie nicht constant sind, selbst dem Systeme $\mathfrak{S}$ angehören, so ist das System $\mathfrak{S}$ endlich.}

Dies ergiebt sich unmittelbar daraus, dass die Grade der Formen $\Phi_1, \Phi_2, \ldots \Phi_\mu$ niedriger sind, als der Grad von $F$. Wendet man also die Darstellung (1) auf die nicht constanten unter den Functionen $\Phi$ an, so gelangt man zu Co\"efficienten von noch niedrigerem Grade und muss also schliesslich bei wiederholter Anwendung dieses Verfahrens auf Constanten kommen.

Daraus folgt nun durch eine sehr einfache Schlussweise, die ich einer mündlichen Mittheilung von \gesperrt{Hurwitz} verdanke, die Endlichkeit des Invariantensystemes $\mathfrak{J}$ einer endlichen Gruppe linearer Substitutionen $S$.
